% !TeX program = xelatex
\documentclass{article}
\usepackage{kotex}
\usepackage[a4paper, left=3cm, right=3cm, top=2.5cm, bottom=2.5cm]{geometry}
\usepackage{amsmath, amssymb}
\usepackage{tcolorbox}
\usepackage{caption}
\usepackage{setspace}
\usepackage{xcolor}
\usepackage{titlesec}
\usepackage{tikz}
\usetikzlibrary{arrows.meta}

\setstretch{1.3}
\renewcommand{\figurename}{그림}
\titleformat{name=\section,numberless}
  {\normalfont\large\bfseries\color{blue}}{}
  {0pt}{}

\title{1.10 일반적인 등거리 변환}
\date{}
\author{}

\begin{document}
\maketitle

\section*{1.10 일반적인 등거리 변환}

\begin{tcolorbox}[
  colback=red!4, colframe=red!60!black,
  title=정리 1.10.1,
  fonttitle=\bfseries, boxrule=0.5mm, arc=3mm]

모든 등거리 변환 $f : \mathbb{E}^n \to \mathbb{E}^n$은
\textbf{회전, 평행 이동 그리고 반사의 합성}이다.

\end{tcolorbox}

\textbf{개략적인 증명} ($n = 2$인 경우; 일반적인 경우도 비슷한 방법으로 증명된다).

$\mathbb{E}^2$에서 좌표를 정하자.
한 쌍의 수직인 수선을 잡고 네 개의 사분면에 번호를 붙인다.
각 점에 한 축($X = x$축)으로부터의 거리 $(\pm y)$와 다른 축($Y = y$축)으로부터의 거리 $(\pm x)$를 재어 좌표 $(x,y)$를 대응시킨다.

등거리 변환 $f$는 $X$와 $Y$를 다른 수직인 직선의 쌍으로 보낸다. 이를 다음과 같이 정하자.
\[
  X' = f(X), \qquad Y' = f(Y)
\]
$T(0,0) = f(0,0)$을 만족하는 평행 이동 $T$를 선택하자.
$f(0,0)$을 중심으로 하고 $T(X)$를 $X'$으로 보내는 회전을 $R$이라 하자
(이와 같은 회전은 두 개가 있는데, 회전각은 $180°$만큼 차이가 난다).
$T(X)$의 양의 부분을 $X'$의 양의 부분으로 보내는 회전을 $R$이라 하자.
그러면 합성 $R \circ T$는 $X$를 $X'$으로 보낸다.
또한 $R \circ T$가 각을 보존하므로 $Y$를 $Y'$으로 보낸다.

\begin{figure}[h]
\centering
\begin{tikzpicture}[>=Stealth, scale=1.0]
  % --- Left: original axes ---
  \draw[->] (-0.3,0) -- (2.2,0) node[right]{$x$};
  \draw[->] (0,-0.3) -- (0,2.2) node[above]{$y$};
  \node[below] at (1.3,0) {$a$};
  \node[left]  at (0,1.1) {$b$};
  \draw[fill=white] (0,1.6) circle (2.5pt);

  % arrow f
  \draw[->, thick] (2.5,1.0) -- (3.8,1.0) node[midway, above]{$f$};

  % --- Right: rotated axes ---
  \begin{scope}[xshift=4.5cm, rotate=35]
    \draw[->] (-0.3,0) -- (2.2,0) node[right]{$x'$};
    \draw[->] (0,-0.3) -- (0,2.2) node[above]{$y'$};
    \node[below] at (1.3,0) {$a$};
    \node[left]  at (0,1.1) {$b$};
    \draw[fill=white] (0,1.6) circle (2.5pt);
  \end{scope}
\end{tikzpicture}
\caption{등거리 변환 $f$에 의한 좌표 축의 이동}
\label{fig:1.51}
\end{figure}

\begin{itemize}\setlength\itemsep{2pt}
  \item \textbf{경우 1}: $R \circ T$가 $Y$의 양의 부분을 $Y'$의 양의 부분으로 보내면,
    $f$와 $R \circ T$가 좌표 축을 동일하게 대응시키므로 $f = R \circ T$.

  \item \textbf{경우 2}: 그렇지 않으면,
    $F$를 $X'$에 대한 반사라 하자.
    그러면 $F \circ R \circ T$는 $X$의 양의 부분을 $X'$의 양의 부분으로 보내고
    $Y$의 양의 부분을 $Y'$의 양의 부분으로 보낸다.
    따라서 $f = F \circ R \circ T$.
\end{itemize}

어느 경우에나 $f$는 평행 이동, 회전, 그리고 반사의 합성이다.
\hfill $\square$

\begin{tcolorbox}[
  colback=red!4, colframe=red!60!black,
  title=따름정리 1.10.1,
  fonttitle=\bfseries, boxrule=0.5mm, arc=3mm]

모든 등거리 변환은 \textbf{반사만의 합성}으로 표현된다.

\end{tcolorbox}

\textbf{증명.}\quad
정리 1.5.1과 문제 1.6.1로부터 $\mathbb{E}^2$에서 모든 회전과 모든 평행 이동은 반사들의 합성이다.
$\mathbb{E}^n$에서도 똑같이 성립한다(증명도 비슷하다).
\hfill $\square$

\bigskip

\begin{tcolorbox}[
  colback=red!4, colframe=red!60!black,
  title=따름정리 1.10.2\quad ($\mathbb{R}^2$에서의 일반적인 등거리 변환 공식),
  fonttitle=\bfseries, boxrule=0.5mm, arc=3mm]

모든 등거리 변환 $f : \mathbb{R}^2 \to \mathbb{R}^2$는 다음과 같이 표현된다.
\[
  f(x, y) \;=\;
  \bigl(a_0 + x\cos\theta \mp y\sin\theta,\;\;
        b_0 + x\sin\theta \pm y\cos\theta\bigr)
  \qquad \text{(부호동순)} \tag{1.6}
\]
여기서 $a_0,\, b_0$와 $\theta$는 상수로서, $\mathbb{R}^2$의 표준 좌표 계에서
$f(0,0) = (a_0, b_0)$이고,
$f$가 \textbf{향을 보존}하면 식 (1.6)에서 \textbf{더하기(+)}, 향을 뒤바꾸면 \textbf{빼기(--)}이다.

\end{tcolorbox}

\textbf{증명.}\quad
$(a_0, b_0) = f(0,0)$이라 하자.
$f(0,0)$을 시점으로 하고, 종점이 $f(1,0)$, $f(0,1)$인 벡터를 각각 $\vec{v}$, $\vec{w}$라 하자.
\[
  \vec{v} = f(1,0) - f(0,0), \qquad \vec{w} = f(0,1) - f(0,0)
\]
$f$가 등거리 변환이므로 $\vec{v}$와 $\vec{w}$는 서로 수직인 단위 벡터이다.
그러므로 다음을 만족하는 각 $\theta$가 존재한다.
\[
  \vec{v} = (\cos\theta,\; \sin\theta), \qquad
  \vec{w} = \pm(-\sin\theta,\; \cos\theta) \tag{1.7}
\]
$\vec{w}$에 대한 식에서 $f$가 향을 보존하면 \textbf{더하기}, 그렇지 않으면 \textbf{빼기}이다.

$f$는 $x$축을 $(a_0, b_0)$를 원점으로 하는 축 $X'$으로 보내고 양의 방향은 $\vec{v}$ 방향에 대응시킨다.
또한 $f$는 $y$축을 $(a_0, b_0)$를 원점으로 하는 축 $Y'$으로 보내고 양의 방향은 $\vec{w}$ 방향에 대응시킨다.
각 점 $f(x,y)$의 $X'Y'$좌표는 $(x,y)$의 $xy$좌표와 동일하므로
\[
  f(x,y) = (a_0,\, b_0) + x\,\vec{v} + y\,\vec{w}
\]
식 (1.7)을 위의 식에 적용하면 결과의 식 (1.6)을 얻는다.
\hfill $\square$

\subsection*{연습문제 1.10}

\begin{enumerate}
  \item \textbf{(문제 1.10.1)}\quad
    $f : \mathbb{E}^2 \to \mathbb{E}^2$가 향을 보존하는 등거리 변환이면
    $f$는 평행 이동이거나 회전임을 보여라.\\
    \textit{(도움말: $f(P) \neq P$인 점 $P$를 선택한다.
    $P$를 원점으로 하고 $f(P)$가 $x$축의 양의 부분에 놓이는 좌표 계를 만들어라.
    선분 $\overline{Pf(P)}$의 수직 이등분선과 선분 $\overline{f(P)f(f(P))}$의 수직 이등분선의
    교점을 $C$라 하자. $f(C) = C$임을 보이고,
    $f$가 $C$를 중심으로 하고 회전각이 $\angle PCf(P)$인 회전임을 유도하여라.
    두 선분의 수직 이등분선이 만나지 않는 경우에는
    $f$가 평행 이동이거나 회전각이 $180°$인 회전임을 보여라.)}

  \item \textbf{(문제 1.10.2)}\quad
    $\mathbb{E}^2$ 또는 $\mathbb{E}^3$에서 모든 등거리 변환의 역함수가 존재함을 보여라.\\
    \textit{(도움말: 문제 1.4.2, 식 (1.1), 식 (1.2)와 함께 정리 1.10.1을 사용하여
    등거리 변환의 역함수를 만들어 보아라.)}
\end{enumerate}

\end{document}
